% ==============================================================================
% Figure 3: Attention as Parallel Transport -- Curvature from Path-Dependence
%
% Shows two paths from x_i to x_k and the curvature gap F_{ijk}.
%
% Usage: % ==============================================================================
% Figure 3: Attention as Parallel Transport -- Curvature from Path-Dependence
%
% Shows two paths from x_i to x_k and the curvature gap F_{ijk}.
%
% Usage: % ==============================================================================
% Figure 3: Attention as Parallel Transport -- Curvature from Path-Dependence
%
% Shows two paths from x_i to x_k and the curvature gap F_{ijk}.
%
% Usage: % ==============================================================================
% Figure 3: Attention as Parallel Transport -- Curvature from Path-Dependence
%
% Shows two paths from x_i to x_k and the curvature gap F_{ijk}.
%
% Usage: \input{figures/fig3_attention_transport.tex}
% Requires: \usepackage{tikz}
%           \usetikzlibrary{arrows.meta, positioning, calc, decorations.pathreplacing}
% ==============================================================================

\begin{figure}[t]
\centering
\begin{tikzpicture}[
    >=Stealth,
    every node/.style={font=\small},
    event/.style={
        circle, draw=black!70, fill=white, line width=0.8pt,
        inner sep=0pt, minimum size=1.0cm, font=\normalsize
    },
    directpath/.style={
        line width=1.4pt, draw=directclr, >=Stealth
    },
    indirectpath/.style={
        line width=1.4pt, draw=indirectclr, >=Stealth
    },
    curvedlabel/.style={
        font=\footnotesize, fill=white, inner sep=2pt, text=#1
    },
]

% === Colors ===
\definecolor{directclr}{RGB}{200, 55, 55}     % red for direct path
\definecolor{indirectclr}{RGB}{41, 98, 165}   % blue for indirect path
\definecolor{curvclr}{RGB}{160, 50, 160}      % purple for curvature
\definecolor{gapclr}{RGB}{160, 50, 160}

% ====================================================================
% EVENT NODES
% ====================================================================
\node[event] (xi) at (0, 0) {$x_i$};
\node[event] (xj) at (4.0, 0) {$x_j$};
\node[event] (xk) at (2.0, 3.8) {$x_k$};

% ====================================================================
% DIRECT PATH: x_i -> x_k (expected transport)
% ====================================================================
\draw[directpath, ->]
    (xi) to[bend left=12]
    node[curvedlabel=directclr, left=3pt, pos=0.5, sloped]
    {$U_{ik}^{\exp}$}
    (xk);

% ====================================================================
% INDIRECT PATH: x_i -> x_j -> x_k (composite transport)
% ====================================================================
\draw[indirectpath, ->]
    (xi) -- node[curvedlabel=indirectclr, below=4pt, pos=0.5]
    {$U_{ij}$}
    (xj);
\draw[indirectpath, ->]
    (xj) to[bend left=12]
    node[curvedlabel=indirectclr, right=3pt, pos=0.5, sloped]
    {$U_{jk}$}
    (xk);

% ====================================================================
% CURVATURE GAP
% ====================================================================
% Shade the region between the two paths to indicate curvature
\fill[curvclr, opacity=0.07]
    (xi) to[bend left=12] (xk) -- (xj) to[bend right=12] (xi) -- cycle;

% Dashed line showing the "gap" between the two transport results at x_k
\coordinate (direct_end) at ($(xi)!0.72!(xk) + (0.15, 0.05)$);
\coordinate (indirect_end) at ($(xj)!0.62!(xk) + (-0.15, 0.05)$);

% Curvature annotation -- placed in the interior of the triangle
\node[
    draw=curvclr, fill=curvclr!6, rounded corners=3pt,
    inner sep=6pt, font=\small, align=center, text width=3.6cm
] (curvlabel) at (2.5, 1.4)
    {$F_{ijk} = U_{ij} \cdot U_{jk}$\\[1pt]$\cdot\, (U_{ik}^{\exp})^{-1}$};

% Arrows from the annotation to the gap region
\draw[->, curvclr, line width=0.7pt, dashed]
    (curvlabel.north west) to[out=140, in=-40]
    ($(xi)!0.60!(xk) + (0.18, 0.0)$);
\draw[->, curvclr, line width=0.7pt, dashed]
    (curvlabel.north east) to[out=40, in=200]
    ($(xj)!0.55!(xk) + (-0.18, 0.0)$);

% ====================================================================
% EVENT LABELS (semantic annotations beneath nodes)
% ====================================================================
\node[font=\scriptsize, text=black!60, below=6pt of xi] {assertion};
\node[font=\scriptsize, text=black!60, below=6pt of xj] {pressure turn};
\node[font=\scriptsize, text=black!60, above=6pt of xk] {response};

% ====================================================================
% PATH LEGEND (bottom)
% ====================================================================
\draw[directpath] (0.3, -1.6) -- (1.2, -1.6);
\node[font=\footnotesize, right=3pt, text=directclr] at (1.2, -1.6)
    {direct (expected) transport \; $U_{ik}^{\exp}$};

\draw[indirectpath] (0.3, -2.1) -- (1.2, -2.1);
\node[font=\footnotesize, right=3pt, text=indirectclr] at (1.2, -2.1)
    {composite transport \; $U_{ij} \cdot U_{jk}$};

\draw[curvclr, dashed, line width=0.7pt] (0.3, -2.6) -- (1.2, -2.6);
\node[font=\footnotesize, right=3pt, text=curvclr] at (1.2, -2.6)
    {holonomy deviation \; $F_{ijk} \neq I_d$ \; (path-dependence)};

% ====================================================================
% CONDITION NOTE (right side)
% ====================================================================
\node[font=\footnotesize, align=left, text width=4.2cm, anchor=north west]
    at (5.0, 3.6)
    {$F_{ijk} = I_d$ \;\textbar\; transport is\\
     path-independent\\[3pt]
     $\|F_{ijk}\|_k > 0$ \;\textbar\; binding\\
     failure in layer $k$};

% ====================================================================
% TRIANGLE EDGE: faint outline to emphasize the loop
% ====================================================================
\draw[black!15, line width=0.5pt]
    (xi) -- (xj) -- (xk) -- cycle;

\end{tikzpicture}
\caption{%
    \textbf{Holonomy deviation as path-dependence of attention transport.}
    Two paths connect event~$x_i$ to event~$x_k$: the direct (expected) transport~$U_{ik}^{\exp}$
    (red), and the composite transport~$U_{ij} \cdot U_{jk}$ through the intermediate
    pressure turn~$x_j$ (blue).
    The holonomy deviation~$F_{ijk} = U_{ij} \cdot U_{jk} \cdot (U_{ik}^{\exp})^{-1}$
    measures the failure of these two paths to agree.
    When~$F_{ijk} = I_d$, transport is path-independent and the standpoint is preserved;
    when~$\|F_{ijk}\|_k > 0$, layer~$k$ has experienced a binding failure.
}
\label{fig:attention-transport}
\end{figure}

% Requires: \usepackage{tikz}
%           \usetikzlibrary{arrows.meta, positioning, calc, decorations.pathreplacing}
% ==============================================================================

\begin{figure}[t]
\centering
\begin{tikzpicture}[
    >=Stealth,
    every node/.style={font=\small},
    event/.style={
        circle, draw=black!70, fill=white, line width=0.8pt,
        inner sep=0pt, minimum size=1.0cm, font=\normalsize
    },
    directpath/.style={
        line width=1.4pt, draw=directclr, >=Stealth
    },
    indirectpath/.style={
        line width=1.4pt, draw=indirectclr, >=Stealth
    },
    curvedlabel/.style={
        font=\footnotesize, fill=white, inner sep=2pt, text=#1
    },
]

% === Colors ===
\definecolor{directclr}{RGB}{200, 55, 55}     % red for direct path
\definecolor{indirectclr}{RGB}{41, 98, 165}   % blue for indirect path
\definecolor{curvclr}{RGB}{160, 50, 160}      % purple for curvature
\definecolor{gapclr}{RGB}{160, 50, 160}

% ====================================================================
% EVENT NODES
% ====================================================================
\node[event] (xi) at (0, 0) {$x_i$};
\node[event] (xj) at (4.0, 0) {$x_j$};
\node[event] (xk) at (2.0, 3.8) {$x_k$};

% ====================================================================
% DIRECT PATH: x_i -> x_k (expected transport)
% ====================================================================
\draw[directpath, ->]
    (xi) to[bend left=12]
    node[curvedlabel=directclr, left=3pt, pos=0.5, sloped]
    {$U_{ik}^{\exp}$}
    (xk);

% ====================================================================
% INDIRECT PATH: x_i -> x_j -> x_k (composite transport)
% ====================================================================
\draw[indirectpath, ->]
    (xi) -- node[curvedlabel=indirectclr, below=4pt, pos=0.5]
    {$U_{ij}$}
    (xj);
\draw[indirectpath, ->]
    (xj) to[bend left=12]
    node[curvedlabel=indirectclr, right=3pt, pos=0.5, sloped]
    {$U_{jk}$}
    (xk);

% ====================================================================
% CURVATURE GAP
% ====================================================================
% Shade the region between the two paths to indicate curvature
\fill[curvclr, opacity=0.07]
    (xi) to[bend left=12] (xk) -- (xj) to[bend right=12] (xi) -- cycle;

% Dashed line showing the "gap" between the two transport results at x_k
\coordinate (direct_end) at ($(xi)!0.72!(xk) + (0.15, 0.05)$);
\coordinate (indirect_end) at ($(xj)!0.62!(xk) + (-0.15, 0.05)$);

% Curvature annotation -- placed in the interior of the triangle
\node[
    draw=curvclr, fill=curvclr!6, rounded corners=3pt,
    inner sep=6pt, font=\small, align=center, text width=3.6cm
] (curvlabel) at (2.5, 1.4)
    {$F_{ijk} = U_{ij} \cdot U_{jk}$\\[1pt]$\cdot\, (U_{ik}^{\exp})^{-1}$};

% Arrows from the annotation to the gap region
\draw[->, curvclr, line width=0.7pt, dashed]
    (curvlabel.north west) to[out=140, in=-40]
    ($(xi)!0.60!(xk) + (0.18, 0.0)$);
\draw[->, curvclr, line width=0.7pt, dashed]
    (curvlabel.north east) to[out=40, in=200]
    ($(xj)!0.55!(xk) + (-0.18, 0.0)$);

% ====================================================================
% EVENT LABELS (semantic annotations beneath nodes)
% ====================================================================
\node[font=\scriptsize, text=black!60, below=6pt of xi] {assertion};
\node[font=\scriptsize, text=black!60, below=6pt of xj] {pressure turn};
\node[font=\scriptsize, text=black!60, above=6pt of xk] {response};

% ====================================================================
% PATH LEGEND (bottom)
% ====================================================================
\draw[directpath] (0.3, -1.6) -- (1.2, -1.6);
\node[font=\footnotesize, right=3pt, text=directclr] at (1.2, -1.6)
    {direct (expected) transport \; $U_{ik}^{\exp}$};

\draw[indirectpath] (0.3, -2.1) -- (1.2, -2.1);
\node[font=\footnotesize, right=3pt, text=indirectclr] at (1.2, -2.1)
    {composite transport \; $U_{ij} \cdot U_{jk}$};

\draw[curvclr, dashed, line width=0.7pt] (0.3, -2.6) -- (1.2, -2.6);
\node[font=\footnotesize, right=3pt, text=curvclr] at (1.2, -2.6)
    {holonomy deviation \; $F_{ijk} \neq I_d$ \; (path-dependence)};

% ====================================================================
% CONDITION NOTE (right side)
% ====================================================================
\node[font=\footnotesize, align=left, text width=4.2cm, anchor=north west]
    at (5.0, 3.6)
    {$F_{ijk} = I_d$ \;\textbar\; transport is\\
     path-independent\\[3pt]
     $\|F_{ijk}\|_k > 0$ \;\textbar\; binding\\
     failure in layer $k$};

% ====================================================================
% TRIANGLE EDGE: faint outline to emphasize the loop
% ====================================================================
\draw[black!15, line width=0.5pt]
    (xi) -- (xj) -- (xk) -- cycle;

\end{tikzpicture}
\caption{%
    \textbf{Holonomy deviation as path-dependence of attention transport.}
    Two paths connect event~$x_i$ to event~$x_k$: the direct (expected) transport~$U_{ik}^{\exp}$
    (red), and the composite transport~$U_{ij} \cdot U_{jk}$ through the intermediate
    pressure turn~$x_j$ (blue).
    The holonomy deviation~$F_{ijk} = U_{ij} \cdot U_{jk} \cdot (U_{ik}^{\exp})^{-1}$
    measures the failure of these two paths to agree.
    When~$F_{ijk} = I_d$, transport is path-independent and the standpoint is preserved;
    when~$\|F_{ijk}\|_k > 0$, layer~$k$ has experienced a binding failure.
}
\label{fig:attention-transport}
\end{figure}

% Requires: \usepackage{tikz}
%           \usetikzlibrary{arrows.meta, positioning, calc, decorations.pathreplacing}
% ==============================================================================

\begin{figure}[t]
\centering
\begin{tikzpicture}[
    >=Stealth,
    every node/.style={font=\small},
    event/.style={
        circle, draw=black!70, fill=white, line width=0.8pt,
        inner sep=0pt, minimum size=1.0cm, font=\normalsize
    },
    directpath/.style={
        line width=1.4pt, draw=directclr, >=Stealth
    },
    indirectpath/.style={
        line width=1.4pt, draw=indirectclr, >=Stealth
    },
    curvedlabel/.style={
        font=\footnotesize, fill=white, inner sep=2pt, text=#1
    },
]

% === Colors ===
\definecolor{directclr}{RGB}{200, 55, 55}     % red for direct path
\definecolor{indirectclr}{RGB}{41, 98, 165}   % blue for indirect path
\definecolor{curvclr}{RGB}{160, 50, 160}      % purple for curvature
\definecolor{gapclr}{RGB}{160, 50, 160}

% ====================================================================
% EVENT NODES
% ====================================================================
\node[event] (xi) at (0, 0) {$x_i$};
\node[event] (xj) at (4.0, 0) {$x_j$};
\node[event] (xk) at (2.0, 3.8) {$x_k$};

% ====================================================================
% DIRECT PATH: x_i -> x_k (expected transport)
% ====================================================================
\draw[directpath, ->]
    (xi) to[bend left=12]
    node[curvedlabel=directclr, left=3pt, pos=0.5, sloped]
    {$U_{ik}^{\exp}$}
    (xk);

% ====================================================================
% INDIRECT PATH: x_i -> x_j -> x_k (composite transport)
% ====================================================================
\draw[indirectpath, ->]
    (xi) -- node[curvedlabel=indirectclr, below=4pt, pos=0.5]
    {$U_{ij}$}
    (xj);
\draw[indirectpath, ->]
    (xj) to[bend left=12]
    node[curvedlabel=indirectclr, right=3pt, pos=0.5, sloped]
    {$U_{jk}$}
    (xk);

% ====================================================================
% CURVATURE GAP
% ====================================================================
% Shade the region between the two paths to indicate curvature
\fill[curvclr, opacity=0.07]
    (xi) to[bend left=12] (xk) -- (xj) to[bend right=12] (xi) -- cycle;

% Dashed line showing the "gap" between the two transport results at x_k
\coordinate (direct_end) at ($(xi)!0.72!(xk) + (0.15, 0.05)$);
\coordinate (indirect_end) at ($(xj)!0.62!(xk) + (-0.15, 0.05)$);

% Curvature annotation -- placed in the interior of the triangle
\node[
    draw=curvclr, fill=curvclr!6, rounded corners=3pt,
    inner sep=6pt, font=\small, align=center, text width=3.6cm
] (curvlabel) at (2.5, 1.4)
    {$F_{ijk} = U_{ij} \cdot U_{jk}$\\[1pt]$\cdot\, (U_{ik}^{\exp})^{-1}$};

% Arrows from the annotation to the gap region
\draw[->, curvclr, line width=0.7pt, dashed]
    (curvlabel.north west) to[out=140, in=-40]
    ($(xi)!0.60!(xk) + (0.18, 0.0)$);
\draw[->, curvclr, line width=0.7pt, dashed]
    (curvlabel.north east) to[out=40, in=200]
    ($(xj)!0.55!(xk) + (-0.18, 0.0)$);

% ====================================================================
% EVENT LABELS (semantic annotations beneath nodes)
% ====================================================================
\node[font=\scriptsize, text=black!60, below=6pt of xi] {assertion};
\node[font=\scriptsize, text=black!60, below=6pt of xj] {pressure turn};
\node[font=\scriptsize, text=black!60, above=6pt of xk] {response};

% ====================================================================
% PATH LEGEND (bottom)
% ====================================================================
\draw[directpath] (0.3, -1.6) -- (1.2, -1.6);
\node[font=\footnotesize, right=3pt, text=directclr] at (1.2, -1.6)
    {direct (expected) transport \; $U_{ik}^{\exp}$};

\draw[indirectpath] (0.3, -2.1) -- (1.2, -2.1);
\node[font=\footnotesize, right=3pt, text=indirectclr] at (1.2, -2.1)
    {composite transport \; $U_{ij} \cdot U_{jk}$};

\draw[curvclr, dashed, line width=0.7pt] (0.3, -2.6) -- (1.2, -2.6);
\node[font=\footnotesize, right=3pt, text=curvclr] at (1.2, -2.6)
    {holonomy deviation \; $F_{ijk} \neq I_d$ \; (path-dependence)};

% ====================================================================
% CONDITION NOTE (right side)
% ====================================================================
\node[font=\footnotesize, align=left, text width=4.2cm, anchor=north west]
    at (5.0, 3.6)
    {$F_{ijk} = I_d$ \;\textbar\; transport is\\
     path-independent\\[3pt]
     $\|F_{ijk}\|_k > 0$ \;\textbar\; binding\\
     failure in layer $k$};

% ====================================================================
% TRIANGLE EDGE: faint outline to emphasize the loop
% ====================================================================
\draw[black!15, line width=0.5pt]
    (xi) -- (xj) -- (xk) -- cycle;

\end{tikzpicture}
\caption{%
    \textbf{Holonomy deviation as path-dependence of attention transport.}
    Two paths connect event~$x_i$ to event~$x_k$: the direct (expected) transport~$U_{ik}^{\exp}$
    (red), and the composite transport~$U_{ij} \cdot U_{jk}$ through the intermediate
    pressure turn~$x_j$ (blue).
    The holonomy deviation~$F_{ijk} = U_{ij} \cdot U_{jk} \cdot (U_{ik}^{\exp})^{-1}$
    measures the failure of these two paths to agree.
    When~$F_{ijk} = I_d$, transport is path-independent and the standpoint is preserved;
    when~$\|F_{ijk}\|_k > 0$, layer~$k$ has experienced a binding failure.
}
\label{fig:attention-transport}
\end{figure}

% Requires: \usepackage{tikz}
%           \usetikzlibrary{arrows.meta, positioning, calc, decorations.pathreplacing}
% ==============================================================================

\begin{figure}[t]
\centering
\begin{tikzpicture}[
    >=Stealth,
    every node/.style={font=\small},
    event/.style={
        circle, draw=black!70, fill=white, line width=0.8pt,
        inner sep=0pt, minimum size=1.0cm, font=\normalsize
    },
    directpath/.style={
        line width=1.4pt, draw=directclr, >=Stealth
    },
    indirectpath/.style={
        line width=1.4pt, draw=indirectclr, >=Stealth
    },
    curvedlabel/.style={
        font=\footnotesize, fill=white, inner sep=2pt, text=#1
    },
]

% === Colors ===
\definecolor{directclr}{RGB}{200, 55, 55}     % red for direct path
\definecolor{indirectclr}{RGB}{41, 98, 165}   % blue for indirect path
\definecolor{curvclr}{RGB}{160, 50, 160}      % purple for curvature
\definecolor{gapclr}{RGB}{160, 50, 160}

% ====================================================================
% EVENT NODES
% ====================================================================
\node[event] (xi) at (0, 0) {$x_i$};
\node[event] (xj) at (4.0, 0) {$x_j$};
\node[event] (xk) at (2.0, 3.8) {$x_k$};

% ====================================================================
% DIRECT PATH: x_i -> x_k (expected transport)
% ====================================================================
\draw[directpath, ->]
    (xi) to[bend left=12]
    node[curvedlabel=directclr, left=3pt, pos=0.5, sloped]
    {$U_{ik}^{\exp}$}
    (xk);

% ====================================================================
% INDIRECT PATH: x_i -> x_j -> x_k (composite transport)
% ====================================================================
\draw[indirectpath, ->]
    (xi) -- node[curvedlabel=indirectclr, below=4pt, pos=0.5]
    {$U_{ij}$}
    (xj);
\draw[indirectpath, ->]
    (xj) to[bend left=12]
    node[curvedlabel=indirectclr, right=3pt, pos=0.5, sloped]
    {$U_{jk}$}
    (xk);

% ====================================================================
% CURVATURE GAP
% ====================================================================
% Shade the region between the two paths to indicate curvature
\fill[curvclr, opacity=0.07]
    (xi) to[bend left=12] (xk) -- (xj) to[bend right=12] (xi) -- cycle;

% Dashed line showing the "gap" between the two transport results at x_k
\coordinate (direct_end) at ($(xi)!0.72!(xk) + (0.15, 0.05)$);
\coordinate (indirect_end) at ($(xj)!0.62!(xk) + (-0.15, 0.05)$);

% Curvature annotation -- placed in the interior of the triangle
\node[
    draw=curvclr, fill=curvclr!6, rounded corners=3pt,
    inner sep=6pt, font=\small, align=center, text width=3.6cm
] (curvlabel) at (2.5, 1.4)
    {$F_{ijk} = U_{ij} \cdot U_{jk}$\\[1pt]$\cdot\, (U_{ik}^{\exp})^{-1}$};

% Arrows from the annotation to the gap region
\draw[->, curvclr, line width=0.7pt, dashed]
    (curvlabel.north west) to[out=140, in=-40]
    ($(xi)!0.60!(xk) + (0.18, 0.0)$);
\draw[->, curvclr, line width=0.7pt, dashed]
    (curvlabel.north east) to[out=40, in=200]
    ($(xj)!0.55!(xk) + (-0.18, 0.0)$);

% ====================================================================
% EVENT LABELS (semantic annotations beneath nodes)
% ====================================================================
\node[font=\scriptsize, text=black!60, below=6pt of xi] {assertion};
\node[font=\scriptsize, text=black!60, below=6pt of xj] {pressure turn};
\node[font=\scriptsize, text=black!60, above=6pt of xk] {response};

% ====================================================================
% PATH LEGEND (bottom)
% ====================================================================
\draw[directpath] (0.3, -1.6) -- (1.2, -1.6);
\node[font=\footnotesize, right=3pt, text=directclr] at (1.2, -1.6)
    {direct (expected) transport \; $U_{ik}^{\exp}$};

\draw[indirectpath] (0.3, -2.1) -- (1.2, -2.1);
\node[font=\footnotesize, right=3pt, text=indirectclr] at (1.2, -2.1)
    {composite transport \; $U_{ij} \cdot U_{jk}$};

\draw[curvclr, dashed, line width=0.7pt] (0.3, -2.6) -- (1.2, -2.6);
\node[font=\footnotesize, right=3pt, text=curvclr] at (1.2, -2.6)
    {holonomy deviation \; $F_{ijk} \neq I_d$ \; (path-dependence)};

% ====================================================================
% CONDITION NOTE (right side)
% ====================================================================
\node[font=\footnotesize, align=left, text width=4.2cm, anchor=north west]
    at (5.0, 3.6)
    {$F_{ijk} = I_d$ \;\textbar\; transport is\\
     path-independent\\[3pt]
     $\|F_{ijk}\|_k > 0$ \;\textbar\; binding\\
     failure in layer $k$};

% ====================================================================
% TRIANGLE EDGE: faint outline to emphasize the loop
% ====================================================================
\draw[black!15, line width=0.5pt]
    (xi) -- (xj) -- (xk) -- cycle;

\end{tikzpicture}
\caption{%
    \textbf{Holonomy deviation as path-dependence of attention transport.}
    Two paths connect event~$x_i$ to event~$x_k$: the direct (expected) transport~$U_{ik}^{\exp}$
    (red), and the composite transport~$U_{ij} \cdot U_{jk}$ through the intermediate
    pressure turn~$x_j$ (blue).
    The holonomy deviation~$F_{ijk} = U_{ij} \cdot U_{jk} \cdot (U_{ik}^{\exp})^{-1}$
    measures the failure of these two paths to agree.
    When~$F_{ijk} = I_d$, transport is path-independent and the standpoint is preserved;
    when~$\|F_{ijk}\|_k > 0$, layer~$k$ has experienced a binding failure.
}
\label{fig:attention-transport}
\end{figure}
